% text-layer.tex — keep the PDF's text layer faithful to the source.
%
% Usable on its own:  pandoc doc.md -H text-layer.tex --pdf-engine=xelatex ...
%
% THE PROBLEM
%
% Many text faces carry a `calt` (contextual alternates) feature that swaps
% minus, hyphen and parentheses for `.case` variants when they sit next to
% capitals or digits — the same marks, optically raised to suit cap height.
%
% Those variants are usually encoded in the Private Use Area. XeTeX maps them
% faithfully, so the PDF's ToUnicode table correctly reports a glyph as, say,
% U+EE6B — and U+EE6B means nothing outside that one font. Extractors then
% disagree: poppler DISCARDS PUA codepoints, so the character silently
% disappears; pypdf keeps them, so you get mojibake where a minus sign belongs.
%
% The result is a PDF that looks perfectly correct and whose text layer is not.
% In a profit-and-loss statement one extractor turns -123,45 into 123,45 and the
% other into a character no parser will read. Neither warns you.
%
% See docs/inter-calt-tounicode.md for the full write-up, a reproducer that
% needs nothing but a stock TeX Live, and the two fixes that do NOT work.
%
% THE FIX
%
% Disable calt on faces used for text that will be extracted. The substitution
% then never happens, the glyph really is U+2212, and no extractor has to
% implement anything optional to agree. \XeTeXgenerateactualtext=1 looks like a
% better fix because it preserves the typography and satisfies poppler — but it
% only adds an OPTIONAL /ActualText hint on top of the unchanged PUA mapping,
% so pypdf and anything else ignoring /ActualText still gets mojibake.
%
% This file APPLIES that fix rather than describing it. \defaultfontfeatures
% sets a default for every font loaded afterwards, so including this fragment is
% enough — you do not have to remember RawFeature={-calt} at each \setmainfont,
% and it still holds when you pass features of your own:
%
%   \setmainfont{YourFace}[Numbers={Proportional,Lining}]   % -calt still applied
%
% An earlier version of this file only explained the defect in a comment. A user
% who included it and then set Inter got the exact bug the file is named after,
% and got no warning, because the LaTeX default font does not exhibit it.
%
% If you deliberately want the case-sensitive punctuation back on a face whose
% text will never be extracted, override locally:
%
%   \setmainfont{YourFace}[RawFeature={+calt}]
%
% The cost is that case-sensitive punctuation. Measured on a 300 dpi A4 render:
% 2,106 of 26 MB of raster bytes change — 0.008% of the page. For a document
% with numbers in it, that is the right trade.

\usepackage{fontspec}
\defaultfontfeatures{RawFeature={-calt}}

% Tabular figures inside long tables so columns of amounts line up.
\usepackage{etoolbox}
\AtBeginEnvironment{longtable}{\addfontfeature{Numbers=Tabular}}
